\Article{11}
The line on the Earth's surface along which $V$ has the same definite value $=V^0$ at every point generally separates the parts of that surface where $V$ is greater than $V^0$ from those where it is smaller. At every point on this line the horizontal magnetic force is evidently perpendicular to it and points toward the side where the larger values of $V$ occur. Let $\mathrm{d}s$ be an infinitely small line in this direction and $V^0+\mathrm{d}V^0$ the value of $V$ at its other endpoint. Then $\frac{\mathrm{d}V^0}{\mathrm{d}s}$ is the intensity of the horizontal magnetic force at that place. Here too, all the points corresponding to the value $V=V^0+\mathrm{d}V^0$ form a second line infinitely close to the first, marking off a \textit{zone} on the Earth's surface in which the values of $V$ lie between $V^0$ and $V^0+\mathrm{d}V^0$, and in which the horizontal intensity is inversely proportional to the width of the zone. Thus, if $V$ is made to vary by infinitely small but equal steps from its smallest value on the Earth's surface to its largest, the entire surface is divided into infinitely many infinitely narrow zones. The horizontal magnetic force is everywhere perpendicular to their dividing lines, and its intensity is inversely proportional to the width of the zones at the places in question. The two extreme values of $V$ then correspond to two points enclosed by the zones, at which the horizontal force becomes $=0$ and \SourceLeaf{138}the total magnetic force can therefore only be vertical. These points are called the Earth's magnetic poles.

The lines dividing the zones are the intersections of the surfaces considered in Article 7 with the Earth's surface. At the poles these surfaces only touch it.

\Article{12}
The form of the system of lines described in the preceding article is really only the simplest type. It might admit a variety of exceptions if every possible distribution of magnetism within the Earth were taken into account. We shall confine ourselves here to a few remarks explaining the exceptions. especially since, given the Earth's \textit{actual} magnetic constitution, the system of lines on its surface does have that form. It certainly admits no large-scale exceptions and may at most have an isolated local exception here and there.

Some physicists have put forward the view that the Earth has two magnetic north poles and two south poles. Yet the most essential condition seems to have been left unmet. No \textit{precise} definition was first given of what was to count as a magnetic pole. We shall use this name for every point on the Earth's surface where the horizontal intensity is $=0$. Generally speaking, the inclination there is therefore $=90^0$. This also includes the singular case, if it should occur, where the total intensity is $=0$. If one wished to call the places where the total intensity attains a maximum magnetic poles, meaning a value greater than at every point in its immediate vicinity, one must keep in mind how different that is from the definition we have given. Such points have no necessary connection with the former ones in either location or number, and using the same name for things of different kinds leads to confusion.

Leaving aside the Earth's actual constitution and taking the question in general, more than two magnetic poles can indeed exist. But it seems to have gone unnoticed that if, for example, there are two north poles, there must also be a third point between them \SourceLeaf{139}that is likewise a magnetic pole, though strictly it is neither a north pole nor a south pole, or, if one prefers, both at once.

The best way to clarify this is to consider our system of lines.

\expandafter\def\csname HistoricalNoteText16\endcsname{Even if taken to be infinitely small, these rings need not be circular and are generally elliptical. The direction of the magnetic needle normal to them therefore coincides with the direction toward $P^*$ at only four places on each ring. One can consequently make substantial errors by taking the intersection of two compass directions extended from considerable distances to be $P^*$ without further inquiry.}
\expandafter\def\csname HistoricalNoteText17\endcsname{We follow the usual terminology here, which gives that name to the point located by Captain Ross, although it is strictly a south pole insofar as the Earth itself is regarded as a magnet.}
If the function $V$ has at a point $P^*$ on the Earth's surface a maximum value $V^*$, with smaller values all around, a series of values decreasing step by step will correspond to a system of rings. Each ring encloses all the preceding ones and the point $P^*$. On every one of these rings, the horizontal magnetic force, or the north pole of the magnetic needle, will point \textit{inward}\HistoricalNote{16}. This is the distinguishing feature of a magnetic north pole{\renewcommand{\thefootnote}{**)}\HistoricalNote{17}}. The rings can clearly be taken so small, or the corresponding values of the function $V$ so close to $V^*$, that any other given point still lies outside.

Let $S$ denote the set of all points on the Earth's surface where the value of $V$ is greater than a given quantity $W$. Clearly, $S$ forms either a single connected region of the surface or several separated regions. \textit{On} the boundary line or lines separating it from the remaining parts, where $V$ is smaller than $W$, we have $V=W$. As $W$ decreases or increases, that region expands or contracts.

Suppose now that $P^{**}$ is a second point of the same kind as $P^*$, so that $V$ has a maximum value $=V^{**}$ there as well. As noted above, we can assign the quantity $W$ a value smaller than $V^*$ and so close to it that $P^{**}$ falls outside the part of $S$ in which $P^*$ lies. Suppose, as we may, that $V^{**}$ is at least as large as $V^*$ and thus also larger than $W$. Then \SourceLeaf{140}$P^{**}$ must also belong to a part of $S$. Therefore $P^*$ and $P^{**}$ both lie in $S$, but in separate parts of $S$.

On the other hand, we can clearly take $W$ so small that $P^*$ and $P^{**}$ lie in one connected part of $S$, since, if we take $W$ small enough, $S$ can cover the whole surface of the Earth.

Now let $W$ pass step by step through all the values from the first to the second. There must be a last value $V^{***}$ for which $P^*,P^{**}$ still lie in separate parts of $S$. As soon as $W$ decreases further, these parts merge into one.

If this merging occurs at a single point $P^{***}$. then the boundary line on which $V=V^{***}$ has the shape of a figure eight that crosses itself at that point, and it is easy to see that the horizontal intensity there must be $=0$. Indeed, that crossing either takes place at a measurable angle or it does not. In the first case, the horizontal force, if it were other than $=0$, would have to be normal to two different tangents, which is absurd. In the second case, the two halves of the figure eight would touch at $P^{***}$, or have the same tangent. The force normal to that tangent could point only toward the interior of one of the two regions enclosed by the figure eight. This is a contradiction, since the value of $V$ increases toward both sides. Thus $P^{***}$ is a true magnetic pole by our definition. It must be regarded as a south pole with respect to the nearest points inside the two loops of the figure eight and as a north pole with respect to the points outside. Figure 1 can illustrate this form of the system of lines.

If the merging occurs at two different places at once, what we have just proved for one point holds for both. It is easy to see that an island-shaped region will then form inside the region enclosing $P^*$ and $P^{**}$. As $W$ continues to decrease, it will grow narrower and must eventually shrink to a true south pole.

Something similar holds if the merging occurs simultaneously at three or more separate points. If it takes place all at once along an entire line, the horizontal force must also vanish at every point of that line.

\SourceLeaf{141}It is also clear that assuming two \textit{south poles} likewise requires the existence of a third pole point, which is neither a south pole nor a north pole, or rather both at once.

\Article{13}
What was developed in the preceding article makes it easy to see. how several conceivable exceptions to the simplest type of our system of lines can arise. All the points corresponding to a given value of $V$ may form a line consisting of several pieces, each of which closes on itself but is completely separate from the others. They may form a line that crosses itself. Finally, they may form a line bordered on both sides by regions of the surface where $V$ is larger than on the line, or by regions where it is smaller on both sides.

We can say that the Earth has no appreciable large-scale departures of this kind from the simplest type. But local departures are quite conceivable where magnetic masses lie just beneath the Earth's surface. At a somewhat considerable distance they no longer have an appreciable effect, but in their immediate vicinity their effect may be strong enough to overwhelm and mask the regularly varying terrestrial magnetic force. In the simplest case, the system of lines in such a region could have the form shown in the $2^{\mathrm{nd}}$ figure.

